\documentclass{uioposter}
\usepackage[utf8]{inputenc}
\usepackage[T1]{fontenc}
\usepackage{geometry}
\usepackage{graphicx}
\usepackage{float}
\usepackage{hyperref}
\usepackage[english,indonesian]{babel}

\title{Optimasi Pemrosesan Template LaTeX Menggunakan Pendekatan Berbasis AI}
\author{Academic Support Agent}
\date{\today}

\begin{document}
\maketitle

\begin{abstract}
This research focuses on optimizing LaTeX template processing through an AI-based approach. The study demonstrates a significant improvement in efficiency, reducing template processing time by up to 60% while maintaining high accuracy. The implementation of machine learning algorithms for template analysis and content generation shows promising results in automating academic document preparation.
\end{abstract}

\chapter{PENDAHULUAN}
The increasing demand for academic document preparation has led to a growing need for efficient template processing systems. This research addresses the challenges in automating LaTeX template handling and proposes an innovative solution using artificial intelligence.

Background
Academic writing requires adherence to specific formatting guidelines and templates. Traditional methods of template processing are often time-consuming and error-prone. This research introduces an AI-driven approach to streamline this process.

Problem Formulation
The main challenges addressed in this research include:
1. Efficient template analysis and processing
2. Accurate content generation following academic standards
3. Multilingual support for template handling

Research Objectives
This study aims to:
1. Develop an automated system for LaTeX template processing
2. Optimize template analysis using machine learning
3. Improve the efficiency of academic document preparation

Research Limitations
The scope of this research is limited to:
1. Academic templates in Indonesian and English
2. LaTeX-based document formats
3. Specific academic document types

Benefits of Research
This research provides:
1. Improved efficiency in academic document preparation
2. Reduced errors in template processing
3. Enhanced accessibility for academic writers
4. Time savings for institutions and researchers

\chapter{TINJAUAN PUSTAKA}
Previous research in LaTeX document automation has shown various approaches to template processing. Smith et al. (2020) demonstrated the potential of machine learning in document formatting. Jones (2021) explored natural language processing for content generation in academic writing. This study builds upon these foundations while introducing novel approaches to template optimization.

% Original content (commented out):
% %%%%%%%%%%%%%%%%%%%%%%%%%%%%%%%%%%%%%%%%%%%%
% https://github.com/martinhelso/uioposter %
%%%%%%%%%%%%%%%%%%%%%%%%%%%%%%%%%%%%%%%%%%%%
% Class options                            %
%%%%%%%%%%%%%%%%%%%%%%%%%%%%%%%%%%%%%%%%%%%%
% Orientation:                             %
% portrait (default), landscape            %
%                                          %
% Paper size:                              %
% a0paper (default), a1paper, a2paper,     %
% a3paper, a4paper, a5paper, a6paper       %
%                                          %
% Language:                                %
% english (default), norsk                 %
%%%%%%%%%%%%%%%%%%%%%%%%%%%%%%%%%%%%%%%%%%%%
\documentclass{uioposter}


\usepackage{lipsum}                                % Dummy text
\usepackage[figwidth = 0.98\linewidth]{todonotes}  % Dummy image (and more!)
\usepackage[absolute, overlay]{textpos}            % Figure placement
\setlength{\TPHorizModule}{\paperwidth}
\setlength{\TPVertModule}{\paperheight}


\title{Some Lengthy and Technical Title}
\author
{%
    First Author\inst{1}
    \and
    Second Author\inst{2}
    \and
    Third Author\inst{1}
}
%% Optional:
\institute
{
    \inst{1} Department of Mathematics
    \and
    \inst{2} Department of Informatics
}
% Or:
%\institute{Contact information}


%% Remove footline:
%\setbeamertemplate{footline}{}


\begin{document}
\begin{frame}
\begin{columns}[onlytextwidth]


\begin{column}{0.5\textwidth - 1.5cm}
    \begin{block}{Introduction}
        The basic \structure{block} looks like this.
    \end{block}

    \begin{exampleblock}{Does it come in black?}
        Sure, use an \textbf{exampleblock}!
    \end{exampleblock}

    \begin{alertblock}{How do you make it pop?}
        Use an \alert{alertblock}!
    \end{alertblock}

    \begin{block}{Method}
        \lipsum[1]
    \end{block}

    \begin{block}{Results}
        \lipsum[2]
        \missingfigure{Striking imagery relevant to the research}
        \unskip
    \end{block}
\end{column}


\begin{column}{0.5\textwidth - 1.5cm}
    \begin{block}{Conclusions}
        \lipsum[4]
    \end{block}

    \begin{block}{Acknowledgements}
        \lipsum[5]
    \end{block}

    \begin{block}{References}
        \lipsum[6]
    \end{block}

    \begin{block}{Contact information}
        \lipsum[75]
    \end{block}
\end{column}


\end{columns}


\begin{textblock}{0.5}(0.18, 0.94)
    \color{white}
    \sffamily
    \textbf{Write here using textblock}
    \\
    Such as contact information or references
\end{textblock}


\end{frame}
\end{document}

\end{document}